
\chapter{2016年北京卷（文）}

\section{单选题}

\begin{problem}
已知集合 \(A=\{x\mid 2<x<4\}\)，\(B=\{x\mid x<3 \text{或} x>5\}\)，则 \(A\cap B=\)（\quad）

\choices
    {\(\{x\mid 2<x<5\}\)}
    {\(\{x\mid x<4 \text{或} x>5\}\)}
    {\(\{x\mid 2<x<3\}\)}
    {\(\{x\mid x<2 \text{或} x>5\}\)}
\end{problem}

\begin{answer}
C
\end{answer}

\begin{solution}
由 \(A=(2,4)\)，\(B=(-\infty,3)\cup(5,+\infty)\)，所以
\[
A\cap B=(2,4)\cap\bigl((-\infty,3)\cup(5,+\infty)\bigr)=(2,3)
\]
故选 C.
\end{solution}

\begin{problem}
复数 \(\dfrac{1+2\i}{2-\i}=\)（\quad）

\choices
    {\(\i\)}
    {\(1+\i\)}
    {\(-\i\)}
    {\(1-\i\)}
\end{problem}

\begin{answer}
A
\end{answer}

\begin{solution}
分子、分母同乘分母的共轭复数 \(2+\i\)，得
\[
\frac{1+2\i}{2-\i}=\frac{(1+2\i)(2+\i)}{(2-\i)(2+\i)}=\frac{2+\i+4\i+2\i^2}{5}=\frac{5\i}{5}=\i
\]
故选 A.
\end{solution}

\begin{problem}
执行如图所示的程序框图，输出 \(s\) 的值为（\quad）

\texfigure[max width=6.0cm]{img_repaint/2016/beijing_liberal/q3_fig1.tex}

\choices
    {\(8\)}
    {\(9\)}
    {\(27\)}
    {\(36\)}
\end{problem}

\begin{answer}
B
\end{answer}

\begin{solution}
程序开始时 \(k=0\)，\(s=0\). 当 \(k\leq2\) 时，执行 \(s=s+k^3\)，再令 \(k=k+1\). 因而循环依次为
\[
(k,s)=(0,0)\longrightarrow(1,0)\longrightarrow(2,1)\longrightarrow(3,9)
\]
此时 \(k=3>2\)，程序停止并输出 \(s=9\). 故选 B.
\end{solution}

\begin{problem}
下列函数中，在区间 \((-1,1)\) 上为减函数的是（\quad）

\choices
    {\(y=\dfrac{1}{1-x}\)}
    {\(y=\cos x\)}
    {\(y=\ln\left( x+1 \right)\)}
    {\(y=2^{-x}\)}
\end{problem}

\begin{answer}
D
\end{answer}

\begin{solution}
在 \((-1,1)\) 上，选项 A 的导数为 \(\dfrac{1}{(1-x)^2}>0\)，选项 C 的导数为 \(\dfrac{1}{x+1}>0\)，二者均为增函数. 选项 B 的导数为 \(-\sin x\)，在 \((-1,0)\) 上为正、在 \((0,1)\) 上为负，因而不是整个区间上的减函数. 选项 D 的导数为 \(\dfrac{\mathrm d}{\mathrm dx}2^{-x}=-(\ln2)2^{-x}<0\)，所以 \(y=2^{-x}\) 在 \((-1,1)\) 上为减函数. 故选 D.
\end{solution}

\begin{problem}
圆 \(\left( x+1 \right)^2+y^2=2\) 的圆心到直线 \(y=x+3\) 的距离为（\quad）

\choices
    {\(1\)}
    {\(2\)}
    {\(\sqrt2\)}
    {\(2\sqrt2\)}
\end{problem}

\begin{answer}
C
\end{answer}

\begin{solution}
圆的圆心为 \(O(-1,0)\). 将直线 \(y=x+3\) 化为 \(x-y+3=0\)，由点到直线的距离公式，得
\[
d=\frac{|(-1)-0+3|}{\sqrt{1^2+(-1)^2}}=\frac{2}{\sqrt2}=\sqrt2
\]
故选 C.
\end{solution}

\begin{problem}
从甲、乙等 \(5\) 名学生中随机选出 \(2\) 人，则甲被选中的概率为（\quad）

\choices
    {\(\dfrac15\)}
    {\(\dfrac25\)}
    {\(\dfrac{8}{25}\)}
    {\(\dfrac{9}{25}\)}
\end{problem}

\begin{answer}
B
\end{answer}

\begin{solution}
从 \(5\) 名学生中任选 \(2\) 人的等可能结果共有 \(\mathrm C_5^2=10\) 个. 若甲被选中，另一人可从其余 \(4\) 人中任选，结果有 \(\mathrm C_4^1=4\) 个. 因此甲被选中的概率为
\[
\frac{\mathrm C_4^1}{\mathrm C_5^2}=\frac{4}{10}=\frac25
\]
故选 B.
\end{solution}

\begin{problem}
已知 \(A(2,5)\)，\(B(4,1)\). 若点 \(P(x,y)\) 在线段 \(AB\) 上，则 \(2x-y\) 的最大值为（\quad）

\choices
    {\(-1\)}
    {\(3\)}
    {\(7\)}
    {\(8\)}
\end{problem}

\begin{answer}
C
\end{answer}

\begin{solution}
设 \(P\) 将线段 \(AB\) 参数化为
\[
P=A+t(B-A)=(2+2t,5-4t)
\]
其中 \(0\leq t\leq1\).
于是
\[
2x-y=2(2+2t)-(5-4t)=-1+8t
\]
当 \(0\leq t\leq1\) 时，\(-1+8t\) 的最大值在 \(t=1\) 时取得，最大值为 \(7\). 故选 C.
\end{solution}

\begin{problem}
某学校运动会的立定跳远和 \(30\) 秒跳绳两个单项比赛分成预赛和决赛两个阶段，表中为 \(10\) 名学生的预赛成绩，其中有三个数据模糊.

\begin{center}
\begin{tabular}{|c|c|c|c|c|c|c|c|c|c|c|}
\hline
学生序号 & \(1\) & \(2\) & \(3\) & \(4\) & \(5\) & \(6\) & \(7\) & \(8\) & \(9\) & \(10\) \\
\hline
立定跳远（单位：米） & \(1.96\) & \(1.92\) & \(1.82\) & \(1.80\) & \(1.78\) & \(1.76\) & \(1.74\) & \(1.72\) & \(1.68\) & \(1.60\) \\
\hline
\(30\) 秒跳绳（单位：次） & \(63\) & \(a\) & \(75\) & \(60\) & \(63\) & \(72\) & \(70\) & \(a-1\) & \(b\) & \(65\) \\
\hline
\end{tabular}
\end{center}

在这 \(10\) 名学生中，进入立定跳远决赛的有 \(8\) 人，同时进入立定跳远决赛和 \(30\) 秒跳绳决赛的有 \(6\) 人，则（\quad）

\choices
    {\(2\)号学生进入 \(30\) 秒跳绳决赛}
    {\(5\)号学生进入 \(30\) 秒跳绳决赛}
    {\(8\)号学生进入 \(30\) 秒跳绳决赛}
    {\(9\)号学生进入 \(30\) 秒跳绳决赛}
\end{problem}

\begin{answer}
B
\end{answer}

\begin{solution}
立定跳远成绩已按从高到低排列，所以进入该项目决赛的是 \(1\) 至 \(8\) 号学生. 由题意，这 \(8\) 人中恰有 \(6\) 人同时进入 \(30\) 秒跳绳决赛.

在这 \(8\) 人中，\(3\)、\(6\)、\(7\) 号的成绩分别为 \(75\)、\(72\)、\(70\) 次. 即使 \(2\) 号和 \(8\) 号的成绩都高于 \(70\) 次，在这三人之前也至多有 \(4\) 人，所以 \(3\)、\(6\)、\(7\) 号均进入 \(30\) 秒跳绳决赛. 还需从 \(1\)、\(2\)、\(4\)、\(5\)、\(8\) 号中选出 \(3\) 人，对应成绩为 \(63\)、\(a\)、\(60\)、\(63\)、\(a-1\).

若 \(5\) 号未进入决赛，则因 \(1\) 号与 \(5\) 号成绩相同，同分学生的决赛资格相同，\(1\) 号也不能进入. 这样剩余应进入的 \(3\) 人只能是 \(2\)、\(4\)、\(8\) 号，其中成绩为 \(60\) 次的 \(4\) 号进入，而成绩为 \(63\) 次的 \(1\)、\(5\) 号均未进入，和按成绩确定决赛资格相矛盾. 因此 \(5\) 号一定进入 \(30\) 秒跳绳决赛，故选 B.
\end{solution}

\section{填空题}

\begin{problem}
已知向量 \(\bs a=(1,\sqrt3)\)，\(\bs b=(\sqrt3,1)\)，则 \(\bs a\) 与 \(\bs b\) 夹角的大小为\fillinblank{}.
\end{problem}

\begin{answer}
\(\dfrac{\pi}{6}\)
\end{answer}

\begin{solution}
由向量数量积公式，\(\bs a\cdot\bs b=1\cdot\sqrt3+\sqrt3\cdot1=2\sqrt3\).
又 \(|\bs a|=\sqrt{1+3}=2\)，\(|\bs b|=\sqrt{3+1}=2\).
设 \(\bs a\) 与 \(\bs b\) 的夹角为 \(\theta\)，则
\[
\cos\theta=\frac{\bs a\cdot\bs b}{|\bs a||\bs b|}=\frac{2\sqrt3}{4}=\frac{\sqrt3}{2}
\]
因为 \(0\leq\theta\leq\pi\)，所以 \(\theta=\dfrac\pi6\).
\end{solution}

\begin{problem}
函数 \(f(x)=\dfrac{x}{x-1}\)（\(x\ge 2\)）的最大值为\fillinblank{}.
\end{problem}

\begin{answer}
\(2\)
\end{answer}

\begin{solution}
在定义域 \(x\geq2\) 上，\(f(x)=\frac{x-1+1}{x-1}=1+\frac1{x-1}\).
由于 \(x-1\geq1\)，所以 \(\dfrac1{x-1}\leq1\)，等号在 \(x=2\) 时取得. 因此 \(f(x)\leq2\)，最大值为 \(2\).
\end{solution}

\begin{problem}
某四棱柱的三视图如图所示，则该四棱柱的体积为\fillinblank{}.

\texfigure[max width=0.72\linewidth]{img_repaint/2016/beijing_liberal/q11_fig1.tex}
\end{problem}

\begin{answer}
\(\dfrac{3}{2}\)
\end{answer}

\begin{solution}
由俯视图可知，四棱柱的底面为梯形，其两条平行边长分别为 \(1\) 和 \(2\)，高为 \(1\)，所以底面积为 \(S_{\text{底}}=\frac{1+2}{2}\times1=\frac32\). 由左视图可知，四棱柱的高为 \(1\). 因而四棱柱的体积为 \(V=S_{\text{底}}\times1=\frac32\).
\end{solution}

\begin{problem}
已知双曲线 \(\dfrac{x^2}{a^2}-\dfrac{y^2}{b^2}=1\)（\(a>0\)，\(b>0\)）的一条渐近线为 \(2x+y=0\)，一个焦点为 \((\sqrt5,0)\)，则 \(a=\fillinblank{}\)，\(b=\fillinblank{}\).
\end{problem}

\begin{answer}
\(1\)，\(2\)
\end{answer}

\begin{solution}
双曲线的渐近线为 \(y=\pm\dfrac ba x\). 已知 \(2x+y=0\)，即 \(y=-2x\)，所以 \(\frac ba=2\). 双曲线的焦点为 \((\pm c,0)\)，且 \(c^2=a^2+b^2\). 由已知焦点得 \(c=\sqrt5\)，故 \(a^2+b^2=5\).
代入 \(b=2a\)，得到 \(5a^2=5\). 因为 \(a>0\)，所以 \(a=1\)，\(b=2\).
\end{solution}

\begin{problem}
在 \(\triangle ABC\) 中，\(\angle A=\dfrac{2\pi}{3}\)，\(a=\sqrt3\,c\)，则 \(\dfrac{b}{c}=\fillinblank{}\).
\end{problem}

\begin{answer}
\(1\)
\end{answer}

\begin{solution}
由余弦定理及 \(\cos\dfrac{2\pi}{3}=-\dfrac12\)，得 \(a^2=b^2+c^2-2bc\cos A=b^2+c^2+bc\).
代入 \(a=\sqrt3\,c\)，得
\(3c^2=b^2+c^2+bc\)，即 \(b^2+bc-2c^2=(b-c)(b+2c)=0\).
由于三角形边长 \(b\)，\(c>0\)，\(b+2c>0\)，所以 \(b=c\)，从而 \(\dfrac bc=1\).
\end{solution}

\begin{problem}
某网店统计了连续三天售出商品的种类情况：第一天售出 \(19\) 种商品，第二天售出 \(13\) 种商品，第三天售出 \(18\) 种商品；前两天都售出的商品有 \(3\) 种，后两天都售出的商品有 \(4\) 种，则该网店

\begin{enumerate}
\item 第一天售出但第二天未售出的商品有\fillinblank{} 种；
\item 这三天售出的商品最少有\fillinblank{} 种.
\end{enumerate}
\end{problem}

\begin{answer}
\(16\)，\(29\)
\end{answer}

\begin{solution}
设第一、二、三天售出的商品集合分别为 \(A\)、\(B\)、\(C\).

第一天售出但第二天未售出的商品数为 \(|A|-|A\cap B|=19-3=16\).

前两天售出的商品共有 \(19+13-3=29\) 种，所以三天售出的商品总数至少为 \(29\). 下面说明这个下界可以取得：令第三天售出的 \(18\) 种商品中有 \(14\) 种属于第一天售出但第二天未售出的商品，另有 \(4\) 种属于第二天和第三天都售出的商品，并且这 \(4\) 种不属于第一天售出商品. 这样既满足后两天共有 \(4\) 种相同商品，又不会使前两天的并集增加. 因此三天售出的商品最少为 \(29\) 种.
\end{solution}

\section{解答题}

\begin{problem}
已知 \(\{a_n\}\) 是等差数列，\(\{b_n\}\) 是等比数列，且 \(b_2=3\)，\(b_3=9\)，\(a_1=b_1\)，\(a_{14}=b_4\).

\begin{enumerate}
\item 求 \(\{a_n\}\) 的通项公式；
\item 设 \(c_n=a_n+b_n\)，求数列 \(\{c_n\}\) 的前 \(n\) 项和.
\end{enumerate}
\end{problem}

\begin{solution}
\begin{enumerate}
\item 设等比数列 \(\{b_n\}\) 的公比为 \(q\). 由 \(b_2=3\)，\(b_3=9\)，得 \(q=\frac{b_3}{b_2}=3\)，所以 \(b_1=\frac{b_2}{q}=1\)，且 \(b_4=b_3q=27\).
设等差数列 \(\{a_n\}\) 的公差为 \(d\). 因为 \(a_1=b_1=1\)，且 \(a_{14}=b_4=27\)，所以
\[
a_{14}=a_1+13d=1+13d=27
\]
解得 \(d=2\). 因此
\[
a_n=a_1+(n-1)d=1+2(n-1)=2n-1
\]

\item 由 \(a_n=2n-1\)，\(b_n=3^{n-1}\)，得
\(c_n=a_n+b_n=2n-1+3^{n-1}\).
于是数列 \(\{c_n\}\) 的前 \(n\) 项和为 \(S_n=\sum_{k=1}^{n}c_k=\sum_{k=1}^{n}(2k-1)+\sum_{k=1}^{n}3^{k-1}\). 由等差数列和等比数列求和公式，\(\sum_{k=1}^{n}(2k-1)=n^2\)，\(\sum_{k=1}^{n}3^{k-1}=\dfrac{3^n-1}{2}\)，所以 \(S_n=n^2+\dfrac{3^n-1}{2}\).
\end{enumerate}
\end{solution}

\begin{problem}
已知函数 \(f(x)=2\sin\omega x\cos\omega x+\cos2\omega x\)（\(\omega>0\)）的最小正周期为 \(\pi\).

\begin{enumerate}
\item 求 \(\omega\) 的值；
\item 求 \(f(x)\) 的单调递增区间.
\end{enumerate}
\end{problem}

\begin{solution}
\begin{enumerate}
\item 由二倍角公式，
\(f(x)=\sin2\omega x+\cos2\omega x=\sqrt2\sin\left(2\omega x+\frac{\pi}{4}\right)\).
因为 \(\omega>0\)，所以 \(f(x)\) 的最小正周期为 \(\dfrac{2\pi}{2\omega}=\dfrac{\pi}{\omega}\). 由题意 \(\dfrac{\pi}{\omega}=\pi\)，解得 \(\omega=1\).

\item 当 \(\omega=1\) 时，
\[
f'(x)=2\cos2x-2\sin2x=2\sqrt2\cos\left(2x+\frac{\pi}{4}\right)
\]
所以 \(f(x)\) 单调递增等价于 \(\cos\left(2x+\frac{\pi}{4}\right)>0\).
即存在 \(k\in\Z\)，使得
\(-\frac{\pi}{2}+2k\pi<2x+\frac{\pi}{4}<\frac{\pi}{2}+2k\pi\).
解得
\(-\frac{3\pi}{8}+k\pi<x<\frac{\pi}{8}+k\pi\). 故 \(f(x)\) 的单调递增区间为 \(\left(-\frac{3\pi}{8}+k\pi,\frac{\pi}{8}+k\pi\right)\)，\(k\in\Z\).
\end{enumerate}
\end{solution}

\begin{problem}
某市居民用水拟实行阶梯水价，每人每月用水量中不超过 \(w\text{ 立方米}\) 的部分按 \(4\text{ 元/立方米}\) 收费，超出 \(w\text{ 立方米}\) 的部分按 \(10\text{ 元/立方米}\) 收费，从该市随机调查了 \(10000\) 位居民，获得了他们某月的用水量数据，整理得到如图频率分布直方图：

\begin{enumerate}
\item 如果 \(w\) 为整数，那么根据此次调查，为使 \(80\%\) 以上居民在该月的用水价格为 \(4\text{ 元/立方米}\)，\(w\) 至少定为多少？
\item 假设同组中的每个数据用该组区间的右端点值代替，当 \(w=3\) 时，估计该市居民该月的人均水费.
\end{enumerate}

\texfigure[max width=0.56\linewidth]{img_repaint/2016/beijing_liberal/q17_fig1.tex}
\end{problem}

\begin{solution}
\begin{enumerate}
\item 由直方图可知，组距为 \(0.5\text{ 立方米}\)，各组的频率依次为 \(0.1\)，\(0.15\)，\(0.2\)，\(0.25\)，\(0.15\)，\(0.05\)，\(0.05\)，\(0.05\). 因而用水量不超过 \(2\text{ 立方米}\) 的居民频率为
\(0.1+0.15+0.2=0.45\).
用水量不超过 \(3\text{ 立方米}\) 的居民频率为
\(0.45+0.25+0.15=0.85\).
因为 \(w\) 为整数，\(w=2\) 时频率只有 \(0.45<0.8\)，而 \(w=3\) 时频率为 \(0.85>0.8\)，所以 \(w\) 至少定为 \(3\).

\item 按题意，用各组区间的右端点代替组内数据时，右端点依次为 \(1\text{ 立方米}\)，\(1.5\text{ 立方米}\)，\(2\text{ 立方米}\)，\(2.5\text{ 立方米}\)，\(3\text{ 立方米}\)，\(3.5\text{ 立方米}\)，\(4\text{ 立方米}\)，\(4.5\text{ 立方米}\)，对应频率依次为 \(0.1\)，\(0.15\)，\(0.2\)，\(0.25\)，\(0.15\)，\(0.05\)，\(0.05\)，\(0.05\). 当 \(w=3\) 时，用水量不超过 \(3\text{ 立方米}\) 的居民每人的水费为用水量的 \(4\) 倍；记用水量为 \(x\text{ 立方米}\)，若 \(x>3\)，则每人的水费为 \(12+10\times\left(x-3\right)\) 元. 所以各组代表用水量对应的水费依次为 \(4\text{ 元}\)，\(6\text{ 元}\)，\(8\text{ 元}\)，\(10\text{ 元}\)，\(12\text{ 元}\)，\(17\text{ 元}\)，\(22\text{ 元}\)，\(27\text{ 元}\)，估计该市居民该月的人均水费为
\[
\bar C=0.1\times4+0.15\times6+0.2\times8+0.25\times10+0.15\times12+0.05\times17+0.05\times22+0.05\times27=10.5\text{ 元}
\]
\end{enumerate}
\end{solution}

\begin{problem}
如图，在四棱锥 \(P-ABCD\) 中，\(PC\perp\) 平面 \(ABCD\)，\(AB\parallel DC\)，\(DC\perp AC\).

\begin{enumerate}
\item 求证：\(DC\perp\) 平面 \(PAC\)；
\item 求证：平面 \(PAB\perp\) 平面 \(PAC\)；
\item 设点 \(E\) 为 \(AB\) 的中点，在棱 \(PB\) 上是否存在点 \(F\)，使得 \(PA\parallel\) 平面 \(CEF\)？说明理由.
\end{enumerate}

\texfigure[max width=0.52\linewidth]{img_repaint/2016/beijing_liberal/q18_fig1.tex}
\end{problem}

\begin{solution}
\begin{enumerate}
\item 因为 \(PC\perp\) 平面 \(ABCD\)，且 \(DC\) 在平面 \(ABCD\) 内，所以 \(DC\perp PC\). 又已知 \(DC\perp AC\). 直线 \(PC\)，\(AC\) 是平面 \(PAC\) 内相交的两条直线，故由直线与平面垂直的判定定理，得
\(DC\perp\) 平面 \(PAC\).

\item 因为 \(AB\parallel DC\)，且 \(DC\perp\) 平面 \(PAC\)，所以 \(AB\perp\) 平面 \(PAC\). 又 \(AB\) 在平面 \(PAB\) 内，故
平面 \(PAB\perp\) 平面 \(PAC\).

\item 取棱 \(PB\) 的中点 \(F\). 因为 \(E\) 是 \(AB\) 的中点，所以在三角形 \(PAB\) 中，中位线定理给出
\(EF\parallel PA\).
又 \(EF\) 在平面 \(CEF\) 内. 由 \(AB\parallel DC\) 及 \(DC\perp AC\)，得 \(AB\perp AC\). 因为 \(E\in AB\)，\(E\ne A\)，且 \(AB\perp AC\)，所以 \(A\notin CE\)：否则 \(A\)，\(C\)，\(E\) 共线，而 \(AE\) 与 \(AC\) 应为同一直线，这与 \(AB\perp AC\) 矛盾. 由于 \(P\notin\) 平面 \(ABCD\)，且 \(F\) 是 \(PB\) 的中点，\(F\notin\) 平面 \(ABCD\)，故平面 \(CEF\) 与平面 \(ABCD\) 的交线为 \(CE\). 因而 \(A\notin\) 平面 \(CEF\)，即 \(PA\) 不在平面 \(CEF\) 内，从而
\(PA\parallel\) 平面 \(CEF\).
这说明存在满足条件的点 \(F\)，且可取 \(F\) 为 \(PB\) 的中点.
\end{enumerate}
\end{solution}

\begin{problem}
已知椭圆 \(C:\dfrac{x^2}{a^2}+\dfrac{y^2}{b^2}=1\) 过点 \(A(2,0)\)，\(B(0,1)\) 两点.

\begin{enumerate}
\item 求椭圆 \(C\) 的方程及离心率；
\item 设 \(P\) 为第三象限内一点且在椭圆 \(C\) 上，直线 \(PA\) 与 \(y\) 轴交于点 \(M\)，直线 \(PB\) 与 \(x\) 轴交于点 \(N\)，求证：四边形 \(ABNM\) 的面积为定值.
\end{enumerate}

\texfigure[max width=0.56\linewidth]{img_repaint/2016/beijing_liberal/q19_fig1.tex}
\end{problem}

\begin{solution}
\begin{enumerate}
\item 将点 \(A(2,0)\)，\(B(0,1)\) 分别代入椭圆方程，得 \(\frac{4}{a^2}=1\)，\(\frac{1}{b^2}=1\). 因为 \(a>0\)，\(b>0\)，所以 \(a=2\)，\(b=1\). 因而椭圆 \(C\) 的方程为 \(\frac{x^2}{4}+y^2=1\). 椭圆的半焦距满足 \(c^2=a^2-b^2=4-1=3\)，所以 \(c=\sqrt3\)，离心率为 \(e=\frac{c}{a}=\frac{\sqrt3}{2}\).

\item 设 \(P(x,y)\) 为第三象限内的点，则 \(x<0\)，\(y<0\)，且
\(x^2+4y^2=4\).
直线 \(PA\) 的参数方程可写为 \(A+\lambda(P-A)\). 令其横坐标为零，得 \(\lambda=\dfrac{2}{2-x}\)，所以 \(M=\left(0,\frac{2y}{2-x}\right)\).
直线 \(PB\) 的参数方程可写为 \(B+\mu(P-B)\). 令其纵坐标为零，得 \(\mu=\dfrac{1}{1-y}\)，所以 \(N=\left(\frac{x}{1-y},0\right)\).

由于 \(x<0\)，\(y<0\)，点 \(N\)，\(M\) 分别位于负半轴上，四边形顶点按 \(A\)，\(B\)，\(N\)，\(M\) 的顺序排列. 由坐标面积公式，记四边形 \(ABNM\) 的面积为 \(S\)，则
\[
S=\frac12\left(2+\frac{2xy}{(1-y)(2-x)}-\frac{x}{1-y}-\frac{4y}{2-x}\right)
\]
将上式因式分解，得
\[
S=\frac12\left(2-\frac{x}{1-y}\right)\left(1-\frac{2y}{2-x}\right)=\frac{(2-x-2y)^2}{2(1-y)(2-x)}
\]
又因为 \(x^2+4y^2=4\)，所以
\[
(2-x-2y)^2-4(1-y)(2-x)=x^2+4y^2-4=0
\]
于是 \((2-x-2y)^2=4(1-y)(2-x)\)，从而 \(S=2\).
故四边形 \(ABNM\) 的面积为定值 \(2\).
\end{enumerate}
\end{solution}

\begin{problem}
设函数 \(f(x)=x^3+ax^2+bx+c\).

\begin{enumerate}
\item 求曲线 \(y=f(x)\) 在点 \((0,f(0))\) 处的切线方程；
\item 设 \(a=b=4\)，若函数 \(f(x)\) 有三个不同零点，求 \(c\) 的取值范围；
\item 求证：\(a^2-3b>0\) 是 \(f(x)\) 有三个不同零点的必要而不充分条件.
\end{enumerate}
\end{problem}

\begin{solution}
\begin{enumerate}
\item 因为 \(f(0)=c\)，且 \(f'(x)=3x^2+2ax+b\)，所以 \(f'(0)=b\). 曲线 \(y=f(x)\) 在点 \((0,f(0))=(0,c)\) 处的切线斜率为 \(b\). 切线方程为 \(y-c=b(x-0)\)，即 \(y=bx+c\).

\item 当 \(a=b=4\) 时，\(f(x)=x^3+4x^2+4x+c\)，且 \(f'(x)=3x^2+8x+4=(3x+2)(x+2)\).
因此 \(f(x)\) 在 \(\left(-\infty,-2\right)\) 和 \(\left(-\dfrac23,+\infty\right)\) 上单调递增，在 \(\left(-2,-\dfrac23\right)\) 上单调递减. 两个极值点处的函数值分别为 \(f(-2)=c\)，\(f\left(-\frac23\right)=c-\frac{32}{27}\).
由于三次函数的最高次项系数为正，要有三个不同零点，必须且只需局部极大值大于零、局部极小值小于零，即 \(c>0\)，\(c-\frac{32}{27}<0\). 所以 \(c\) 的取值范围为 \(0<c<\frac{32}{27}\).

\item 设 \(f(x)\) 有三个不同的实数零点 \(x_1\)，\(x_2\)，\(x_3\). 由韦达定理，\(a=-(x_1+x_2+x_3)\)，\(b=x_1x_2+x_2x_3+x_3x_1\).
于是
\[
\begin{aligned}
a^2-3b
&=(x_1+x_2+x_3)^2-3(x_1x_2+x_2x_3+x_3x_1)\\
&=\frac12\left((x_1-x_2)^2+(x_2-x_3)^2+(x_3-x_1)^2\right)>0.
\end{aligned}
\]
因此 \(a^2-3b>0\) 是 \(f(x)\) 有三个不同零点的必要条件.

但该条件不是充分条件. 例如取 \(a=b=4\)，\(c=2\)，则 \(a^2-3b=16-12=4>0\).
然而此时两个极值点处的函数值为 \(f(-2)=2>0\)，\(f\left(-\frac23\right)=2-\frac{32}{27}=\frac{22}{27}>0\).
函数图象的局部极大值和局部极小值都在 \(x\) 轴上方，所以函数只有一个实数零点，不具有三个不同零点. 故 \(a^2-3b>0\) 不是充分条件.
\end{enumerate}
\end{solution}
